\uuid{c49I}
\titre{Vecteurs}
\auteur{bodin ; hanani ; mzari}
\organisation{exo7}
\licence{CC BY-NC-SA 4.0}
\datecreate{2020-01-20}
\niveau{L1}
\module{Géométrie}
\chapitre{Géométrie du plan}
\sousChapitre{Distance, norme, produit scalaire}
\difficulte{1}
\theme{vecteurs}
\exo7qcmid{lille-1.geometrie-plane.002}

\contenu{
\begin{question}
\qclassification{140.01,140.02}
\qtags{motcle=vecteurs}
On considère les vecteurs $\vec{u}=(2,-1)$ et $\vec{v}=(1,-4)$. Quelles sont les bonnes réponses ?
\begin{answers}
    \bad{La norme de $\vec{u}$ est $\|\vec{u}\|=2-1=1$.}
    \good{La norme de $\vec{u}$ est $\|\vec{u}\|=\sqrt{5}$.}
    \bad{Le produit scalaire de $\vec{u}$ et $\vec{v}$ est $\vec{u}\cdot \vec{v}=(2-1)+(1-4)=-3$.}
    \good{Le produit scalaire de $\vec{u}$ et $\vec{v}$ est $\vec{u}\cdot \vec{v}=6$.}
\end{answers}
\begin{explanations}
Penser aux définitions : si $\vec{u}=(x,y)$ et $\vec{v}=(x',y')$ alors
$\vec{u}\cdot \vec{v} = xx'+yy'$ et $\|\vec{u}\| = \sqrt{\vec{u}\cdot \vec{u}}
= \sqrt{x^2+y^2}$.
\end{explanations}
\end{question}
}
